%----------------------------------------------------------------------------
\chapter{Related Work and Gap Analysis}\label{chap:related}
%----------------------------------------------------------------------------
\todo[inline]{About 3--4 pages. Criteria B/2 (literature, 0--4) and D/1 (originality, 0--5). Compare approaches along explicit criteria and end with the gap your work fills. Cite research papers and current industrial solutions, not only websites.}

\section{Comparison criteria}
TODO Name the criteria, derived from the problem in \cref{chap:intro}, and say why each matters.

\section{Existing approaches}
TODO Group the related work into families and discuss each family against the criteria.

\begin{table}[!ht]
  \centering
  \caption{Comparison of related approaches. \checkmark{} supported, $\circ$ partially supported, -- not supported. TODO Fill in.}
  \label{tab:related}
  \begin{tabularx}{\linewidth}{X c c c c}
    \toprule
    \textbf{Approach} & \textbf{Criterion 1} & \textbf{Criterion 2} & \textbf{Criterion 3} & \textbf{Open source} \\
    \midrule
    TODO Approach A \cite{lamport1978time} & \checkmark & $\circ$ & -- & \checkmark \\
    TODO Approach B \cite{w3c2025vcdm}    & -- & \checkmark & -- & -- \\
    \midrule
    \textbf{This work}                     & \checkmark & \checkmark & \checkmark & \checkmark \\
    \bottomrule
  \end{tabularx}
\end{table}

\section{Gap}
TODO One paragraph that states the gap in one sentence, based on \cref{tab:related}. This paragraph is the novelty argument.
